\documentclass[11pt]{scrartcl}
\usepackage[a4paper,margin=2.6cm]{geometry}
\usepackage{../common/rlstudy}

\title{Extended reading --- Week 11: Contextual Bandits and Off-Policy Evaluation}
\date{\today}
\begin{document}
\maketitle

\section*{Five papers to read this week}

The five papers below cover the spine of the bandit + OPE literature you
will lean on for synthetic insurance / ATS work: the original UCB1 analysis, the LinUCB
news-recommendation paper that founded contextual bandits as a deployable
technique, the LinTS analysis that turned posterior sampling into a
defensible default, the Doubly Robust estimator (the only OPE tool you
should reach for first), and BanditNet (counterfactual risk minimisation
at deep-net scale).

\subsection*{1. Auer, Cesa-Bianchi, Fischer --- \emph{Finite-time Analysis of the Multiarmed Bandit Problem} (Machine Learning, 2002)}

The paper that turned the 30-year-old Lai--Robbins asymptotic into a
finite-time, instance-dependent regret bound. UCB1 is presented in two
pages; the rest is a careful Hoeffding-based proof and an extension to
the variance-aware UCB1-tuned.

\textbf{Why it matters this week.} Every bandit paper since cites this one.
The proof technique --- ``either a bad concentration event happened, or
$N_a \geq 8 \log T / \Delta_a^2$'' --- is the canonical bandit proof
template. You will see it repeated for LinUCB, KL-UCB, and MOSS.

\begin{quote}
``At each play, UCB1 maintains a confidence bound for the mean reward of
each arm. The policy chooses the arm with the largest upper confidence
bound. We show that the expected regret of this policy after any number
$n$ of plays is at most $O(\log n)$.''
\end{quote}

\emph{Commentary.} The single sentence above is the entire premise of OFU
(optimism in the face of uncertainty). The $O(\log n)$ rate is tight up to
constants; the instance-dependent $1/\Delta_a$ scaling is the bit you must
keep in mind when reasoning about \emph{when} bandit algorithms are slow ---
small $\Delta_a$ means hard problem, and no algorithm can escape this.

\subsection*{2. Li, Chu, Langford, Schapire --- \emph{A Contextual-Bandit Approach to Personalized News Article Recommendation} (WWW 2010, arXiv:1003.0146)}

The paper that put contextual bandits on the industry map. Two algorithms
(LinUCB-disjoint, LinUCB-hybrid), deployed on Yahoo!\ Today Module,
evaluated by a clever \emph{replay} estimator that turns logged uniform
exposures into unbiased offline evaluations.

\textbf{Why it matters this week.} The disjoint-arms LinUCB you implement
in \texttt{project/agents.py} is the LinUCB from this paper, with the
tightened confidence radius from Abbasi-Yadkori et al.\ 2011. The hybrid
variant (a shared feature across arms) is what you would deploy in synthetic insurance
when most brokers have few observations.

\begin{quote}
``If we constrain the design of the experiment so that articles are chosen
uniformly at random in the recorded events, then replay yields an unbiased
offline evaluation.''
\end{quote}

\emph{Commentary.} This is the simplest OPE estimator possible --- IPS
with a uniform behaviour policy collapses to ``keep matches, drop misses''.
It works because uniform exposure gives perfect overlap. Production loggers
are almost never uniform; you need the IPS / DR machinery from this week.

\subsection*{3. Agrawal \& Goyal --- \emph{Thompson Sampling for Contextual Bandits with Linear Payoffs} (ICML 2013, arXiv:1209.3352)}

The frequentist regret analysis of LinTS. Earlier work (Chapelle \& Li
2011) had shown empirically that TS dominates UCB on industrial datasets;
this paper gives the theoretical scaffolding.

\textbf{Why it matters this week.} The proof structure --- bound the
probability that a sub-optimal arm's posterior dominates the optimal arm's
posterior --- generalises directly to Beta-TS (Agrawal \& Goyal 2012),
NeuralTS, and meta-Thompson. The frequentist bound $\tilde O(d^{3/2} \sqrt{T})$
is loose by a factor of $\sqrt{d}$ compared to LinUCB; Russo \& Van Roy 2014
recovered the $d \sqrt T$ rate via an information-theoretic argument.

\begin{quote}
``Despite its potential and empirical success, theoretical understanding
of Thompson Sampling has been quite limited. \dots\ We prove a high
probability regret bound for Thompson Sampling for contextual bandits
with linear payoff functions of $\tilde O(d^{3/2} \sqrt{T})$.''
\end{quote}

\emph{Commentary.} The empirical-vs-theoretical gap is a recurring theme.
Take the message: TS works, the theory is conservative, the production
default is ``LinTS with $v$ tuned by replay on a held-out log''.

\subsection*{4. Dud\'ik, Langford, Li --- \emph{Doubly Robust Policy Evaluation and Learning} (ICML 2011, arXiv:1103.4601)}

The DR estimator's modern home. Combines the IPS-style importance-weighted
correction with a Direct-Method baseline, in exactly the form we derived
in chapter 2.

\textbf{Why it matters this week.} DR is the OPE tool you should default
to. Lower variance than IPS, lower bias than the Direct Method, and the
formal double-robustness guarantee makes it the safest choice when you
are unsure whether the propensities or the reward model are correct ---
and in practice you are always unsure.

\begin{quote}
``The estimator is doubly robust: the estimator is consistent when either
the reward model or the propensity model is correct. The doubly robust
estimator typically improves upon the direct method and the IPS estimator
since it uses both reward and propensity information.''
\end{quote}

\emph{Commentary.} The same DR construction underlies the modern
``self-normalised DR'' (Thomas \& Brunskill 2016), ``MAGIC'' for sequential
RL (Thomas \& Brunskill 2016), and the SWITCH-DR thresholding (Wang,
Agarwal, Dud\'ik 2017). One technique, three decades of papers.

\subsection*{5. Joachims, Swaminathan, de Rijke --- \emph{Deep Learning with Logged Bandit Feedback} (ICLR 2018; BanditNet, arXiv:1707.07724)}

Turns the counterfactual risk minimisation (CRM) principle of Swaminathan
\& Joachims 2015 into a deep-net training loss. The key trick is the
\emph{equivariant} SNIPS objective with a Lagrangian rewrite that makes
SGD stable on logged bandit data.

\textbf{Why it matters this week.} CRM + BanditNet is the only published
recipe that makes ``train a deep policy directly against logged data''
work without falling into propensity overfitting. You will use the same
loss in Week 12 to train the ATS re-ranker against logged shortlists.

\begin{quote}
``[\dots] direct application of standard supervised learning techniques
to logged bandit feedback is hindered by the bias inherent in propensity
weighting and by the high variance of the IPS estimator. We address these
issues by introducing a counterfactual risk minimisation principle and
deriving a deep-network training objective.''
\end{quote}

\emph{Commentary.} The Lagrangian rewrite to handle the self-normalisation
divisor is non-obvious; read the appendix carefully. The same numerical
recipe works for ranking (CLTR), recommendation, and ad allocation.

\section*{Bonus paper: Saito \& Joachims --- \emph{Off-Policy Evaluation for Large Action Spaces via Embeddings} (ICML 2022, arXiv:2202.06317)}

The most important applied OPE advance of the last five years. Standard
IPS variance scales with the number of arms $K$; when actions admit a
natural embedding $e(a) \in \Real^{d_e}$ with $d_e \ll K$, the
Marginalised IPS (MIPS) estimator replaces per-action propensities by
marginal embedding propensities and cuts variance from $O(K/n)$ to
$O(d_e / n)$.

\textbf{Why it matters.} synthetic insurance may have 6 brokers today but 600 next
year; ATS has 10\,000+ active requisitions at any time. Without MIPS, OPE
on these scales is statistically void.

\begin{quote}
``We show that the variance of the proposed estimator depends on the
number of effective action embeddings, rather than the number of actions
$|\mathcal{A}|$.''
\end{quote}

\emph{Commentary.} The required assumption --- conditional independence
of the reward from $a$ given $(x, e)$ --- is the embedding-sufficiency
assumption. For broker matching, a good embedding is the broker's
historical conversion profile by lead category; for jobs, the requisition
skill vector. Practical: pick the most expressive embedding that still
satisfies sufficiency by audit (residual check on a holdout).

\section*{Bonus paper: Bottou et al.\ --- \emph{Counterfactual Reasoning and Learning Systems} (JMLR 2013)}

The Microsoft adCenter paper. 81 pages. Reads like a textbook on how to
turn a real ad-auction system into a counterfactual-evaluation laboratory.
Co-authored by Leon Bottou (SGD), John Langford (vowpal wabbit, contextual
bandits at Microsoft), and the adCenter ML team.

\textbf{Why it matters.} This is the most honest published account of
running OPE in production. It covers: how to randomise enough to log
propensities without burning revenue; how to detect propensity overlap
failures; how to compose multiple counterfactual estimators; how to write
the dashboard a non-ML stakeholder will actually trust.

\begin{quote}
``Causal inference is no harder than predictive analytics when the proper
counterfactual data has been collected. Designing the data collection is
the hard part.''
\end{quote}

\emph{Commentary.} The line above should be tattooed on every applied
ML engineer's wrist. The data collection (logging propensities, version
hashes, sensitive attributes, delay outcomes) is 80\% of the project; the
estimator is 20\%. Read this paper twice before touching a synthetic insurance /
ATS production logger.

\section*{Bonus paper: Bietti, Agarwal, Langford --- \emph{A Contextual Bandit Bake-off} (JMLR 2021, arXiv:1802.04064)}

A 47-page systematic empirical comparison of 21 contextual-bandit
algorithms across 500+ datasets. The most useful paper if you actually
have to pick an algorithm for production.

\textbf{Why it matters.} It collapses three decades of bandit research
into a small number of practical recommendations. Spoiler: ``cover with
softmax exploration on top of boosted trees, with an offline DR-corrected
loss'' wins on the median problem; LinUCB / LinTS are competitive on
linear-friendly datasets but dominated when reward is non-linear in the
features.

\begin{quote}
``We find that an algorithm that runs cover-greedy with an online softmax
exploration policy, combined with a doubly-robust loss for offline
optimisation, outperforms more sophisticated algorithms across the suite
of datasets we consider.''
\end{quote}

\emph{Commentary.} The lesson is not ``cover always wins''. The lesson is
that the algorithm you pick matters less than (i) whether your features
are good, (ii) whether you log propensities cleanly, and (iii) whether
your offline evaluator is doubly robust. Defaults beat sophistication.

\section*{Glossary of terms used in this week's papers}

\begin{description}\setlength\itemsep{2pt}
  \item[Behaviour / logging policy ($\pi_b$).] The policy that generated
  the offline log. You evaluate counterfactuals against $\pi_b$.
  \item[Target policy ($\pi$).] The policy whose value $V(\pi)$ you want
  to estimate without deploying it.
  \item[Positivity / overlap.] $\pi_b(a \mid x) > 0$ wherever
  $\pi(a \mid x) > 0$. Required for IPS / DR unbiasedness.
  \item[Importance ratio ($w_i$).] $\pi(a_i \mid x_i) / \pi_b(a_i \mid x_i)$.
  Large weights cause variance; capping is the standard variance-control trick.
  \item[ESS.] Effective sample size,
  $(\sum w_i)^2 / \sum w_i^2$. The size of the equivalent i.i.d.\ sample.
  \item[CRM.] Counterfactual risk minimisation (Swaminathan \& Joachims
  2015): train a policy by minimising a self-normalised IPS objective with
  variance regularisation.
  \item[DR.] Doubly robust: combine IPS-weighted residual correction with
  a direct-method baseline; unbiased if either component is correct.
  \item[MIPS.] Marginalised IPS (Saito \& Joachims 2022); use embedding
  propensities instead of action propensities.
  \item[Propensity overfitting.] Re-fitting $\pi_b$ from logs leaks reward
  into the propensity; biases IPS. Log $\pi_b$ at decision time.
  \item[Delayed feedback.] Reward arrives after a stochastic delay.
  Biases recent-policy evaluation downward. Mitigate by censoring or
  survival modelling (Chapelle 2014; Vernade et al.\ 2017).
\end{description}

\section*{Where to go next}

\begin{itemize}
  \item \textbf{Textbook}: Lattimore \& Szepesv\'ari, \emph{Bandit
  Algorithms} (Cambridge UP, 2020) --- chapters 7--10 cover UCB, KL-UCB,
  MOSS; 19--22 cover LinUCB and LinTS; 28--32 cover OPE and CRM.
  \item \textbf{Empirical state of the art}: Bietti, Agarwal, Langford,
  \emph{A Contextual Bandit Bake-off} (JMLR 2021) --- the right reference
  for which algorithm to use on real logs (spoiler: bag-of-trees with
  bootstrap-TS and an offline doubly-robust correction).
  \item \textbf{Large action spaces}: Saito \& Joachims, \emph{Off-Policy
  Evaluation for Large Action Spaces via Embeddings} (ICML 2022; MIPS).
  The variance-from-$K$-to-$d_e$ trick is the most important practical
  advance of the past five years for matching applications.
  \item \textbf{Production}: Bottou et al., \emph{Counterfactual Reasoning
  and Learning Systems} (JMLR 2013) --- the Microsoft adCenter paper that
  spells out the full propensity-logging, OPE, and counterfactual-training
  pipeline. Mandatory reading before any production deployment.
\end{itemize}

\emph{Sourced from} \texttt{research/bandits\_matching\_sota.md} sections
1--3 (stochastic, contextual, and OPE) and section 12 (production case
studies).

\section*{Annotated reading order}

If you have only one weekend, read in this order:
\begin{enumerate}
  \item Bottou et al.\ 2013 (the production map). Skim sections 1--3, deep-read 4--6.
  \item Auer, Cesa-Bianchi, Fischer 2002 (UCB1). Two pages of algorithm, 12 of proof.
  \item Li et al.\ 2010 (LinUCB). Algorithm, replay evaluation, production results.
  \item Dud\'ik, Langford, Li 2011 (DR). Three pages of derivation.
  \item Bietti, Agarwal, Langford 2021 (Bake-off). Skim the result tables, read sections 2 and 5.
\end{enumerate}

If you have one focused day after the weekend, add:
\begin{enumerate}\setcounter{enumi}{5}
  \item Agrawal \& Goyal 2013 (LinTS). The frequentist analysis, focus on Lemma 1.
  \item Joachims, Swaminathan, de Rijke 2018 (BanditNet). The CRM Lagrangian.
  \item Saito \& Joachims 2022 (MIPS). The embedding-sufficiency assumption.
\end{enumerate}

If you are about to deploy this in production, also read:
\begin{itemize}
  \item Wu, Srikant et al.\ 2016 (Conservative Bandits).
  \item Chapelle 2014 (Delayed Feedback).
  \item Lattimore \& Szepesv\'ari 2020, chapter 32 (off-policy evaluation chapter).
\end{itemize}

\end{document}
